% ------------------------------------------------------------------
%  Capitolo 19 — Dalla scuola all'università
%  Parte V
%  Ricerca di riferimento: ricerca 04 §5–6; 02 §9
%  Voci bibliografiche principali: mcdaniel2020, biwer2020, biwer2023, hartwig2012, dignath2008, yang2021
%  Epigrafe: ✔ verificata sul testo (vedi ricerca 03 e registro, ricerca 00)
%  Scheda del capitolo: ricerca/06_indice-ragionato.md, capitolo 19
% ------------------------------------------------------------------
\chapter{Dalla scuola all'università}
\label{cap:transizione}

\epigrafe[Il giorno che segue è allievo del giorno che precede.]{Discipulus est prioris posterior dies.}{Publilio Siro, Sententiae}

\iniziale{Q}{uesto capitolo} \segnaposto{apertura: da scrivere — vedi la scheda nell'indice ragionato}.

\section{Il salto dell'autonomia}

\segnaposto{da scrivere}

\section{Che cosa la scuola non ha insegnato}

\segnaposto{da scrivere}

\section{Il primo anno: perché si abbandona}

\segnaposto{da scrivere}

\section{Imparare a imparare: un corso per le matricole}

\segnaposto{da scrivere}

\section{Ai maturandi, e ai loro insegnanti}

\segnaposto{da scrivere}

\begin{daricordare}
\segnaposto{il principio del capitolo in due righe}
\end{daricordare}

\begin{domande}
  \item \segnaposto{domanda di recupero su questo capitolo}
  \item \segnaposto{domanda di applicazione a una materia che studi}
  \item \segnaposto{domanda di ripresa su un capitolo precedente}
\end{domande}

\begin{sintesi}
  \item \segnaposto{punto di sintesi}
\end{sintesi}

\finecapitolo
